\chapter{การเรียนรู้เชิงลึกและโมเดลโครงข่ายน้ำหนักเบาสำหรับการตรวจวัดมลพิษทางอากาศ}
\label{chap:ch04}

\section*{แผนบริหารการสอนประจำบทที่ 4}
\addcontentsline{toc}{section}{แผนบริหารการสอนประจำบทที่ 4}

\noindent\textbf{หัวข้อเนื้อหาประจำบท}
\begin{enumerate}
    \item สถาปัตยกรรมโครงข่ายประสาทเทียมแบบคอนโวลูชันน้ำหนักเบา
    \item โครงข่ายการถดถอยเชิงพื้นที่และเวลาและการผสมผสานข้อมูลอุตุนิยมวิทยา
    \item เกณฑ์การประเมินประสิทธิภาพและผลการทดสอบโมเดล
\end{enumerate}

\noindent\textbf{วัตถุประสงค์เชิงพฤติกรรม}
\begin{enumerate}
    \item อธิบายหลักการ ทฤษฎีฟิสิกส์ และแบบจำลองทางสิ่งแวดล้อมประจำบทที่ 4 ได้อย่างถูกต้อง
    \item วิเคราะห์ สังเคราะห์ และประยุกต์ใช้เครื่องมือตรวจวัดหรือกระบวนการแปรรูปพลังงานได้อย่างมีประสิทธิภาพ
    \item คำนวณ ประเมินผลทางเทอร์โมไดนามิกส์ ทัศนศาสตร์ หรือการลดก๊าซเรือนกระจกได้อย่างถูกต้องตามหลักวิชาการ
\end{enumerate}

\noindent\textbf{การวัดและประเมินผล}
\begin{table}[H]
\centering\small
\begin{tabularx}{\textwidth}{p{0.55\textwidth} p{0.25\textwidth} c}
\toprule
\textbf{เกณฑ์การประเมินผล} & \textbf{เครื่องมือวัดผล} & \textbf{สัดส่วนคะแนน} \\
\midrule
1. ทักษะการปฏิบัติการทดลองและการคำนวณ & แบบประเมินทักษะห้องแล็บ & 40\% \\
2. รายงานผลการทดลองและการวิเคราะห์ & การตรวจดาต้าชีทและรายงาน & 30\% \\
3. แบบฝึกหัดทบทวนและคำถามท้ายบท & แบบทดสอบท้ายบทเรียน & 20\% \\
4. จิตพิสัยและการมีส่วนร่วม & การเข้าเรียนและการตรงต่อเวลา & 10\% \\
\midrule
\textbf{รวมทั้งสิ้น} & & \textbf{100\%} \\
\bottomrule
\end{tabularx}
\end{table}
\newpage

% ==========================================================
% เนื้อหาบทเรียนประจำบทที่ 4
% ==========================================================

\section{สถาปัตยกรรมโครงข่ายประสาทเทียมแบบคอนโวลูชันน้ำหนักเบา}
\index{สถาปัตยกรรมโครงข่ายประสาทเทียมแบบคอนโวลูชันน้ำหนักเบา}

การศึกษาและทำความเข้าใจเกี่ยวกับสถาปัตยกรรมโครงข่ายประสาทเทียมแบบคอนโวลูชันน้ำหนักเบา (Lightweight CNN Architectures for Mobile Devices) มีความสำคัญอย่างยิ่งในงานฟิสิกส์สิ่งแวดล้อมและพลังงาน โดยมีรายละเอียดสาระสำคัญดังนี้

การประมวลผลโมเดล Deep Learning บนอุปกรณ์สมาร์ทโฟนมีข้อจำกัดด้านพลังงานประมวลผล หน่วยความจำ และแบตเตอรี่ จึงจำเป็นต้องใช้สถาปัตยกรรมโครงข่ายน้ำหนักเบา (Lightweight Neural Networks):

1. MobileNetV3: ใช้โครงสร้าง Depthwise Separable Convolution ผสานกับ Inverted Residual Block และกลไก Squeeze-and-Excitation (SE) ร่วมกับการค้นหาสถาปัตยกรรมด้วยระบบอัตโนมัติ (Neural Architecture Search; NAS) ทำให้ลดจำนวนพารามิเตอร์ลงเหลือไม่ถึง 4 ล้านตัว
2. EfficientNet-Lite: ใช้การปรับสเกลโครงข่ายแบบผสมผสาน (Compound Scaling) ทั้งความลึก ความกว้าง และความละเอียดของภาพอย่างสมดุล
โมเดลเหล่านี้สามารถแปลงเป็นรูปแบบ TensorFlow Lite (TFLite) หรือ ONNX Runtime เพื่อรันการทำนายค่าฝุ่น PM2.5 ภายในเวลาเสี้ยววินาทีบนชิปประมวลผล NPU ของสมาร์ทโฟน

\section{โครงข่ายการถดถอยเชิงพื้นที่และเวลาและการผสมผสานข้อมูลอุตุนิยมวิทยา}
\index{โครงข่ายการถดถอยเชิงพื้นที่และเวลาและการผสมผสานข้อมูลอุตุนิยมวิทยา}

การศึกษาและทำความเข้าใจเกี่ยวกับโครงข่ายการถดถอยเชิงพื้นที่และเวลาและการผสมผสานข้อมูลอุตุนิยมวิทยา (Spatio-Temporal and Multimodal Fusion) มีความสำคัญอย่างยิ่งในงานฟิสิกส์สิ่งแวดล้อมและพลังงาน โดยมีรายละเอียดสาระสำคัญดังนี้

ความแม่นยำของการประเมิน PM2.5 จากภาพถ่ายจะเพิ่มขึ้นอย่างมีนัยสำคัญเมื่อนำ ข้อมูลทางอุตุนิยมวิทยา (Meteorological Metadata) เข้ามาร่วมประมวลผลแบบมัลติโมดัล (Multimodal Fusion):
- ข้อมูลภาพ (Visual Features): สกัดผ่านชั้น Convolutional layers เพื่อดึงข้อมูลความเบลอ สี และคอนทราสต์
- ข้อมูลอุตุนิยมวิทยา (Tabular Features): ความชื้นสัมพัทธ์ (RH), อุณหภูมิ, ความเร็วลม, และเวลาของวัน (Solar Hour)
- โครงสร้างการเชื่อมต่อข้อมูล (Feature Concatenation): นำเวกเตอร์คุณลักษณะของภาพมาต่อรวมกับเวกเตอร์ข้อมูลสภาพอากาศผ่านชั้น Dense layers ก่อนส่งเข้าสู่ชั้นเอาต์พุตเพื่อทำนายค่าความเข้มข้น PM2.5 ($\mu g/m^3$)

\section{เกณฑ์การประเมินประสิทธิภาพและผลการทดสอบโมเดล}
\index{เกณฑ์การประเมินประสิทธิภาพและผลการทดสอบโมเดล}

การศึกษาและทำความเข้าใจเกี่ยวกับเกณฑ์การประเมินประสิทธิภาพและผลการทดสอบโมเดล (Model Evaluation Metrics and Benchmark Results) มีความสำคัญอย่างยิ่งในงานฟิสิกส์สิ่งแวดล้อมและพลังงาน โดยมีรายละเอียดสาระสำคัญดังนี้

ดัชนีชี้วัดทางสถิติที่ใช้ประเมินความแม่นยำของโมเดลทำนาย PM2.5:
- ค่าสัมประสิทธิ์การตัดสินใจ ($R^2$): ควรมีค่ามากกว่า 0.85 บ่งบอกว่าโมเดลอธิบายความแปรปรวนของข้อมูลได้ดี
- ค่ารากที่สองของความคลาดเคลื่อนเฉลี่ย (Root Mean Square Error; RMSE):
\[ \text{RMSE} = \sqrt{\frac{1}{N} \sum\_{i=1}^N (y\_i - \hat{y}\_i)^2} \]
- ค่าความคลาดเคลื่อนสัมบูรณ์เฉลี่ย (Mean Absolute Error; MAE):
\[ \text{MAE} = \frac{1}{N} \sum\_{i=1}^N |y\_i - \hat{y}\_i| \]
ในการทดสอบด้วยภาพถ่ายเขตเมือง โมเดล MobileNetV3 แบบ Multimodal สามารถบรรลุค่า MAE ต่ำกว่า 6.5 $\mu g/m^3$ เทียบกับสถานีตรวจวัดมาตรฐาน

\section{ขั้นตอนและวิธีปฏิบัติการทดลอง}
\index{ขั้นตอนการทดลองประจำบทที่ 4}

\subsection{การแปลงโมเดล Deep Learning เป็น TFLite และการ Quantization แบบ 8 บิต}
\begin{conceptbox}[การแปลงโมเดล Deep Learning เป็น TFLite และการ Quantization แบบ 8 บิต]
1. เทรนโมเดล MobileNetV3 บนเฟรมเวิร์ก PyTorch หรือ TensorFlow จนกระทั่งได้ค่า Loss ต่ำสุด
2. บันทึกโมเดลหลักในรูปแบบ SavedModel
3. ใช้อรรถประโยชน์ TFLiteConverter เพื่อแปลงโมเดล
4. กำหนดกลยุทธ์การลดขนาดโมเดลแบบ Post-Training Quantization (INT8): แปลงค่าน้ำหนักจาก Float32 เป็น Int8
5. ตรวจสอบขนาดไฟล์โมเดล: ขนาดลดลงจาก 16 MB เหลือประมาณ 4.2 MB โดยความแม่นยำลดลงไม่เกิน 1.5\%
\end{conceptbox}

\subsection{การทดสอบรันแอปพลิเคชันตรวจวัด PM2.5 บนระบบปฏิบัติการ Android}
\begin{conceptbox}[การทดสอบรันแอปพลิเคชันตรวจวัด PM2.5 บนระบบปฏิบัติการ Android]
1. สร้างโปรเจกต์ Android Studio นำเข้าโมเดลไฟล์ `.tflite` ลงในโฟลเดอร์ assets
2. เขียนโค้ด CameraX เพื่อเปิดกล้องและจับภาพเฟรมพรีวิว
3. ทำการ Crop ภาพส่วนของขอบฟ้าและอาคารเข้าสู่ขนาดอินพุต $224 \times 224$ พิกเซล
4. ดึงค่าพิกัด GPS อุณหภูมิ และความชื้นจากเซ็นเซอร์ภายในเครื่อง
5. รันคำสั่ง `interpreter.run()` แสดงผลลัพธ์ค่า PM2.5 และระดับสีแถบคุณภาพอากาศ (AQI) บนหน้าจอทันที
\end{conceptbox}

\section{ตารางบันทึกผลการทดลองและการอภิปรายผล}
\begin{table}[H]
\centering\small
\caption{ตารางบันทึกผลการทดลองและการตรวจวัดพารามิเตอร์ประจำบทที่ 4}
\label{tab:ch04_datasheet}
\begin{tabularx}{\textwidth}{p{0.3\textwidth} p{0.3\textwidth} X}
\toprule
\textbf{พารามิเตอร์ที่ตรวจวัด} & \textbf{ค่าที่บันทึกได้} & \textbf{การแปลผลและการวิเคราะห์ทางฟิสิกส์} \\
\midrule
การทดสอบชุดที่ 1 & ค่าอยู่ในเกณฑ์มาตรฐาน & สอดคล้องกับแบบจำลองทฤษฎี \\
การทดสอบชุดที่ 2 & ค่ามีแนวโน้มเพิ่มขึ้นตามสัดส่วน & ยืนยันสมมติฐานการวิจัยเชิงประจักษ์ \\
ประสิทธิภาพรวม & ร้อยละ 88.5 & แสดงผลสัมฤทธิ์ในระดับยอมรับได้ \\
\bottomrule
\end{tabularx}
\end{table}

\section{คำถามทบทวนและแบบฝึกหัดท้ายบท}
\begin{enumerate}
    \item จงอธิบายข้อดีของการใช้ Depthwise Separable Convolution ใน MobileNet เมื่อเทียบกับ Standard Convolution
    \item เหตุใดการป้อนข้อมูลความชื้นสัมพัทธ์ (Relative Humidity) เข้าร่วมกับภาพถ่ายจึงช่วยเพิ่มความแม่นยำในการทำนายค่า PM2.5 ได้อย่างมาก
    \item หากโมเดลให้ค่า $R^2 = 0.88$ และ $\text{RMSE} = 5.8 \ \mu g/m^3$ จงแปลความหมายเชิงสถิติต่อการนำไปใช้งานจริง
\end{enumerate}
